\documentclass[10pt]{article}

\usepackage[utf8]{inputenc}

\usepackage[T1]{fontenc}

\usepackage{amsmath}

\usepackage{amsfonts}

\usepackage{amssymb}

\usepackage[version=4]{mhchem}

\usepackage{stmaryrd}

\usepackage{bbold}



\begin{document}

[6] инвариант автоморфизмов конечной $W^{*}$-алгебры (см. Неймана аләебра), обобщающий энтропию метрич. динамич. системы. Исследована [7] энтропия автоморфизмов произвольной $W^{*}$-алгебры относительно состояния $\varphi$.



Jum.: [1] C o n n e s A., "Ann. sci. École norm. supér.», 1975, t. 8, p. 383-419; [2] Г о л о д е у В. Я., "Успехи матем. наук", 1978, т. 33 , в. 1, с. 43-94; [3] A r a k i H., "Lecture Notes in Math.», 1978, № 650, p. 66-84; [4] С и н а й Я. Г., А н ш е л е в и ч B. В., "Успехи матем. наук», 1976, т. 31, B. 4, c. 151-67; [5] L a n c e E. C., "Invent. math.", 1976, t. 37, p. 201-14; [6] C o n n e s A., S t $\phi$ r m e r E., "Acta math.", 1975, v. 134, p. $289-306$; [7] С т е п и н А. М., Ӹ у х о в А. Г.,

«Изв. ВУЗов. Математика", 1982, № 8, с. $52-60$. М. Степин.

А. М.



ЭРГОДИЧЕСКОЕ МНОЖЕСТВО в фазовом пространстве $X$ (являющемся метризуемым компактом) топологической динамической системь $\left\{S_{t}\right\}$ (потока или каскада), отвечающее нормированной әргодической инвариантной мере $\mu,-$ множество точек $x \in X$, для к-рых:



a) при любой непрерывной функции $f: X \rightarrow \mathbb{R}$ «временно́е среднее»



\[

\frac{1}{T} \int_{0}^{T} f\left(S_{t} x\right) d t \longrightarrow \int_{X} f d \mu \text { при } T \longrightarrow \infty

\]



б) $\mu(U)>0$ для любой окрестности $U$ точки $x$.



Точка, для к-рой при любой непрерывной $f$ существует предел временно́го среднего в а), наз. к в а з и р ег. уля р н о й. Для такой точки әтот предел имеет вид $\int f d \mu$, где $\mu-$ нек-рая нормированная инвариантная мера, однако не обязательно эргодическая. Если при этом выполняется б), то точка наз. т о ч ко й п л о т н о с т и, а если $\mu$ әргодическая, то точка наз. т р а н з и т и в н о й; при выполнении обоих этих условий она наз. р е г у л я р н о й. Множество нерегулярных точек имеет меру нуль относительно любой инвариантной нормированной меры. Разбиение множества регулярных точек на Э. м., отвечающие различным $\mu$, представляет собой наиболее сильную реализацию идеи о разложении динамич. системы на әргодич. компоненты (требующую и наиболее сильных ограничений на систему).



Э. м. введены Н. Н. Боголюбовым и Н. М. Крыловым [1], другие изложения и обсуждения различных обобщений и смежных вопросов см. в лит. при статьях Инвариантная мера, 1) и Метрическая транзитивHocmb.



Лит.: [1] Б ог о лю б о в Н. Н., Избр. труды, т. 1, К., 1969, с. 411-63. Д. В. Аносов.



ӘРГОДИЧНОСТЬ ди н а ми ч е с ко й с и с т ем ы - свойство, расскатриваемое в эргодической теории. Первоначально оно определялось для каскадов $\left\{T^{k}\right\}$ и потоков $\left\{T_{t}\right\}$ с конечной инвариантной мерой $\mu$ следующим образом: если на фазовом пространстве $W$ задана функция $f \in L_{1}(W, \mu)$, то для почти каждой точки $w$ существует временно́е среднее вдоль траектории әтой точки, т. е.



или



\[

\lim _{n \rightarrow \infty} \frac{1}{n} \sum_{k=0}^{n-1} f\left(T^{k} w\right)

\]



\[

\left.\lim _{T \rightarrow \infty} \frac{1}{T} \int_{0}^{T} f\left(T_{t} w\right) d t\right),

\]



к-рое совпадает с пространственным средним (т. е. с $\left.\frac{1}{\mu(W)} \int f d \mu\right)$. В этом случае говорят также об ә р г од и ч н о с т и меры $\mu$. В частности, для любого измеримого множества $A \subset W$ среднее время пребывания в $A$ траектории почти каждої точки пропорционально $\mu(A)$ (на самом деле это свойство эквивалентно Ә.). Когда была доказана Биркгодја эргодическая теорема, стало ясно, тто Ј. эквивалентна метрическод транзи- тиєности. В связи с этим стали говорить об Ә., понимая под ней метрич. транзитивность, в более общей ситуации, когда уже не приходится говорить о совпадении временны́х и пространственных средних (системы с бесконечной инвариантной или квазиинвариантной мерой, не только потоки и каскады, но и более общие группы п полугруппы преобразований).



Д. В. Аносов



ЭРДЕІІА ЗАДАЧА - задача 0 существовании в $n$-мерном евклидовом пространстве $E_{n}$ множества точек, каждые три из к-рых образуют нетупоугөльный треугольник (с в о й с т в о Э р д е II a), содержащего более чем $2^{n}$ элементов. Поставлена П. Әрдешом (см. [1]); им же было высказано предположение (доказапное в [2]), что задача имеет отрицательный ответ и если множество, обладающее свойством Эрдеша, содержит $2^{n}$ элементов, то это может быть в том и только в том случае, когда оно исчерпывает множество вершин $n$-мерного прямоугольного параллелепипеда из $E_{n}$. Доказательство этого утверждения явилось решением и т. н. 3 а д а ч и К л и: чему равно число вершин $m(K)$ многогранника $K \subset E_{n}$, каждые две вершины к-рого лежат в различных параллельных опорных гиперплоскостях к многограннику $K$ (с в о й с т в о К л и). Если множество $N \subset E^{n}$ обладает свойством Эрдеша, то выпуклая оболочка $\operatorname{conv} N$ множества $N$ представляет собой многогранник $M=\operatorname{conv} N$, обладающий свойством Кли, и $m(M)$ равно мощности множества $N$. Если многогранник $K$ обладает свойством Кли, то $m(K) \leqslant 2^{n}$. Равенство $m(K)=2^{n}$ характеризует $n$-мерный параллелепипед в множестве всех многогранников, обладающих свойством Кли.



Э. 3. связана с Xадвигера гипотезой: $b(M)=m(M)$. Jum.: [1] E r d ö s P., "Michigan Math. J.", 1957, No 4, p. $291-300 ;$; 2] D a n z e r 'L., G r ü n b a u m B., "Math. Z.", ЭРЛАНГА РАСПРЕДЕЛЕНИЕ, ә р л а н $\mathrm{r}$ о в с $\boldsymbol{\kappa}$ о е р а с п р е д е л е н и е, - сосредоточенное на $(0, \infty)$ распределение вероятностей $\mathrm{c}$ плотностью



\[

p(x)=\frac{(n \mu)^{n}}{\Gamma(n)} x^{n-1} e^{-n \mu x}, x>0,

\]



где целое $n \geqslant 1$ и действительное $\mu>0$ - параметры. Характеристич. функция Э. р. имеет вид



\[

\left(1-\frac{i t}{n \mu}\right)^{-n}

\]



а математич. ожидание и дисперсия - соответственно $1 / \mu$ и $1 / n \mu^{2}$.



Э. р. есть специальный случай гамма-распределения: $p(x)=\alpha g_{\lambda}(\alpha x)$, где $g_{\lambda}(y)$ есть плотность гамма-распределения при $\lambda=n$ и где $\alpha=n \mu$. При $n=1$ Э. р. совпадает с показательным распределением с параметром $\mu$. Э. р. с параметрами $n$ и $\mu$ может быть найдено как распределение суммы $n$ независимых случайных величин, цмеющих одинаковое показательное распределение с параметром $n \mu$. При $n \rightarrow \infty$ Э. р. стремится к вырожденному распределению, сосредоточенному в точке $1 / \mu$.



Выделение Э. р. из системы гамма-распределений объясняется использованием Э. р. в теории массового обслуживания. Во многих случайных процессах массового обслуживания Ә. р. появляется как распределение промежутков между случайными событиями или как распределение времен обслуживания. Иногда Э. р. наз. гамма-распределение с плотностью



\[

\frac{\alpha^{n}}{\Gamma(n)} x^{n-1} e^{-\alpha x}, \quad x>0 .

\]



Названо по имени А. Эрланга (А. Erlang), построившего первые математич. модели в задаqах массового обслуживания.



A. B. Прохоров.



JIит.: [1] С а а т и Г. ЈI., Элементы теории массового обслуживания ил ее приложения, пер. с англ., 2 изд., М., 1971.





\end{document}